\documentclass[11pt]{article}
\usepackage[letterpaper,margin=1in]{geometry}
\usepackage[T1]{fontenc}
\usepackage{lmodern}
\usepackage{microtype}
\usepackage{amsmath,amssymb,amsthm,mathtools}
\usepackage{xcolor}
\usepackage{hyperref}
\hypersetup{colorlinks=true,linkcolor=blue!45!black,
 citecolor=blue!45!black,urlcolor=blue!45!black,
 pdftitle={Further quantitative bounds for bounded-batch discovery},
 pdfauthor={Christopher D. Long}}
\usepackage{fancyhdr}
\pagestyle{fancy}
\fancyhf{}
\fancyhead[L]{\small Further bounds for batch discovery}
\fancyhead[R]{\small Research draft}
\fancyfoot[C]{\thepage}
\setlength{\headheight}{14pt}
\setlength{\emergencystretch}{2em}
\numberwithin{equation}{section}
\newtheorem{theorem}{Theorem}[section]
\newtheorem{lemma}[theorem]{Lemma}
\newtheorem{proposition}[theorem]{Proposition}
\newtheorem{corollary}[theorem]{Corollary}
\theoremstyle{remark}
\newtheorem{remark}[theorem]{Remark}
\newcommand{\E}{\mathbb E}
\newcommand{\Pp}{\mathbb P}
\newcommand{\R}{\mathbb R}
\newcommand{\C}{\mathbb C}
\newcommand{\ind}{\mathbf 1}
\newcommand{\Var}{\operatorname{Var}}
\newcommand{\Bin}{\operatorname{Bin}}
\newcommand{\dK}{d_{\mathrm K}}
\newcommand{\Normal}{\mathcal N(0,1)}
\newcommand{\Phibar}{\overline\Phi}
\title{Further quantitative bounds for bounded-batch discovery}
\author{Christopher D. Long}
\date{November 26, 2026\\[3pt]\small Research draft}

\begin{document}
\maketitle

\begin{abstract}
We apply the variance-linear cumulant bounds for Poissonized bounded-batch
discovery to relative tail approximation. Write
$S=(X_h-\E X_h)/\sqrt V$ and $\delta=W/\sqrt V$, where each batch has total
nonnegative species weight at most $W$. Uniformly over the batch law, time,
countable species space, and overlap pattern, we obtain a saddlepoint
approximation and a full Cram\'er--Petrov correction with relative error
$O(\delta(1+x))$ for $0\le x\le c/\delta$. In a fixed smaller range,
$\log\{\Pp(S\ge x)/\overline\Phi(x)\}=O(\delta(1+x^3))$, while subtracting
$\kappa_3(S)x^3/6$ leaves an error
$O(\delta(1+x)+\delta^2x^4)$. Thus the Gaussian relative approximation holds
uniformly for $x=o(\delta^{-1/3})$, and the cubic-corrected approximation
holds for $x=o(\delta^{-1/2})$. Poissonized singleton examples show that both
exponents are optimal for the model class. All estimates cover both signs
and both strict and non-strict tails, without a lattice assumption. The
cumulant transfer is classical Cram\'er--Petrov theory; we give a
self-contained exponential-tilt proof with explicit conservative constants
for this application. We also prove uniform $O(W)$ mean and $O(W^2)$
variance corrections after a deterministic number of batches, with the
exact residual variance $V_h(n)-nf_h(n)^2$; the remainder orders are sharp.
For multiple coverage, an exact integrated Poisson envelope replaces a
threshold factor of order $r$ by one of order $\sqrt r$, improving the
bracket bounds, generic normal approximation, and actual-variance
concentration. These last two advances do not establish the corresponding
optimal Berry--Esseen rates.
\end{abstract}

\section{Model, imported cumulants, and main conclusions}\label{sec:main}

Let independent, identically distributed random finite batches $A_1,A_2,\ldots$ from a countable
species set arrive according to an independent rate-one Poisson process
$N(t)$. Assign species weights $h_i\ge0$ and assume
\begin{equation}\label{eq:batchweight}
 \sum_{i\in A_1}h_i\le W\quad\text{almost surely},\qquad W>0.
\end{equation}
At a fixed time $t$, put
\[
 X_h=\sum_{i\in\bigcup_{k=1}^{N(t)}A_k}h_i,
 \qquad V=\Var(X_h)>0,
 \qquad S=\frac{X_h-\E X_h}{\sqrt V},
 \qquad \delta=\frac W{\sqrt V}.
\]
Every batch is finite, but their cardinalities need not have a uniform bound.
Unweighted discovery with $|A_1|\le B$ is the special case $h_i=1$ and $W=B$.
The deterministic bound $X_h\le WN(t)$ gives all exponential moments.

Sections~\ref{sec:main}--\ref{sec:sharpness} develop the relative-tail
conclusions. Section~\ref{sec:uniform-correction} separately treats
fixed-sample mean and variance corrections, and
Section~\ref{sec:sharper-threshold} treats multiple coverage. Each later
section states its own notation and proof inputs.

We import the cumulant estimate and zero-free moment generating function
from \cite[Theorem~1.2 and Proposition~3.2]{LongOptimal}:
\begin{equation}\label{eq:importedcumulants}
 |\kappa_k(S)|\le2k!(4\delta)^{k-2},\qquad k\ge3,
\end{equation}
and $\E e^{zS}$ is zero-free for $|z|<1/(2\delta)$. These results already
include countably many species and arbitrary overlap. The relative-tail argument uses
them as explicit inputs. No assertion about the distribution after a
deterministic number of batches is needed for that argument.

Write $\Phi$ and $\phi$ for the standard normal distribution function and
density, and $\Phibar(x)=1-\Phi(x)$. To state all endpoint conventions at
once, let $\epsilon\in\{-1,1\}$ and
$\diamond\in\{\ge,>\}$, and set
\[
 p_{\epsilon,\diamond}(x)=\Pp(\epsilon S\diamond x),\qquad x\ge0.
\]
For $\epsilon=-1$, these are respectively $\Pp(S\le-x)$ and
$\Pp(S<-x)$. Let
\begin{equation}\label{eq:Ksign}
 K_\epsilon(z)=\log\E e^{z\epsilon S},\qquad K_\epsilon(0)=0,
\end{equation}
using the analytic logarithm near zero. On the real line it is the ordinary
real logarithm of a positive moment generating function.

\begin{theorem}[Saddlepoint and full Cram\'er correction]\label{thm:tails}
For each sign $\epsilon$ and
\begin{equation}\label{eq:mainrange}
 0\le x\le\frac1{2048\delta},
\end{equation}
there is a unique $h=h_\epsilon(x)\in[0,(1024\delta)^{-1}]$ satisfying
$K_\epsilon'(h)=x$. Define
\begin{align}
 v_\epsilon(x)&=K_\epsilon''(h),&
 a_\epsilon(x)&=h\sqrt{v_\epsilon(x)},\notag\\
 I_\epsilon(x)&=hx-K_\epsilon(h),&
 \Lambda_\epsilon(x)&=\frac{x^2}{2}-I_\epsilon(x),
 \label{eq:saddlequantities}
\end{align}
and, for $a\ge0$, put
\begin{equation}\label{eq:Jdef}
 J(a)=e^{a^2/2}\Phibar(a).
\end{equation}
For either choice of $\diamond$,
\begin{equation}\label{eq:saddlebound}
 \left|\frac{p_{\epsilon,\diamond}(x)}
 {e^{-I_\epsilon(x)}J(a_\epsilon(x))}-1\right|
 \le1800\delta(1+x),
\end{equation}
and
\begin{equation}\label{eq:fullcorrection}
 \left|\frac{p_{\epsilon,\diamond}(x)}
 {\Phibar(x)e^{\Lambda_\epsilon(x)}}-1\right|
 \le2500\delta(1+x).
\end{equation}
The correction obeys
\begin{align}
 |\Lambda_\epsilon(x)|&\le10\delta x^3,
 \label{eq:Lambdabound}\\
 \left|\Lambda_\epsilon(x)
       -\frac{\epsilon\kappa_3(S)x^3}{6}\right|
 &\le500\delta^2x^4.
 \label{eq:cubicLambda}
\end{align}
All constants are universal. No smallness assumption on $\delta$ is needed
for these statements, although the relative error bounds are useful when
$\delta(1+x)$ is small.
\end{theorem}

The function $I_\epsilon$ in \eqref{eq:saddlequantities} is the Legendre
transform of the real cumulant generating function at $x$: strict convexity
and $K_\epsilon'(h)=x$ identify its maximizer. The correction
$\Lambda_\epsilon=x^2/2-I_\epsilon$ retains all orders of the Cram\'er--Petrov
series in this range.

\begin{corollary}[Logarithmic relative-error bounds]\label{cor:logbounds}
Suppose
\begin{equation}\label{eq:smallrange}
 0<\delta\le\frac1{16384},\qquad
 0\le x\le\frac1{16384\delta}.
\end{equation}
Then, for both signs and both endpoint conventions,
\begin{align}
 \left|\log\frac{p_{\epsilon,\diamond}(x)}
 {\Phibar(x)e^{\Lambda_\epsilon(x)}}\right|
 &\le5000\delta(1+x),
 \label{eq:logfull}\\
 \left|\log\frac{p_{\epsilon,\diamond}(x)}{\Phibar(x)}\right|
 &\le12000\delta(1+x^3),
 \label{eq:loggaussian}\\
 \left|\log\frac{p_{\epsilon,\diamond}(x)}{\Phibar(x)}
       -\frac{\epsilon\kappa_3(S)x^3}{6}\right|
 &\le5000\delta(1+x)+500\delta^2x^4.
 \label{eq:logcubic}
\end{align}
\end{corollary}

Consequently, for any sequence of models with $\delta\to0$, the following
statements are uniform over $0\le x\le b_\delta$, both signs, and both
endpoint conventions:
\begin{align}
 \frac{p_{\epsilon,\diamond}(x)}{\Phibar(x)}&=1+o(1)
 &&\text{if }b_\delta=o(\delta^{-1/3}),
 \label{eq:gaussianwindow}\\
 \frac{p_{\epsilon,\diamond}(x)}
 {\Phibar(x)\exp\{\epsilon\kappa_3(S)x^3/6\}}&=1+o(1)
 &&\text{if }b_\delta=o(\delta^{-1/2}),
 \label{eq:cubicwindow}\\
 \frac{p_{\epsilon,\diamond}(x)}
 {\Phibar(x)e^{\Lambda_\epsilon(x)}}&=1+o(1)
 &&\text{if }b_\delta=o(\delta^{-1}).
 \label{eq:fullwindow}
\end{align}
Section~\ref{sec:sharpness} proves that the exponents $1/3$ and $1/2$ in
the first two statements are optimal for this model class. The numerical
constants in the finite bounds are deliberately conservative.

\subsection*{Classical provenance}

The transfer from cumulants to corrected relative tails is classical.
D\"oring, Jansen, and Schubert~\cite[Theorem~2.1]{DJS} give a
Cram\'er--Petrov approximation under
$|\kappa_k(S)|\le(k-2)!\Delta^{-(k-2)}$. Our input satisfies that
condition with
\begin{equation}\label{eq:Delta}
 \Delta=(48\delta)^{-1},
\end{equation}
because
\[
 2k!(4\delta)^{k-2}
 \le(k-2)!(48\delta)^{k-2},\qquad k\ge3.
\]
Indeed, $2k(k-1)\le12^{k-2}$ holds with equality at $k=3$ and follows
thereafter by induction. Their theorem therefore already supplies the
universal order of the full-correction error and the Gaussian relative-tail
range. The results here apply that theory to the discovery cumulants, make
the endpoint conventions explicit, and provide a self-contained tilt proof
of the displayed constants and cubic remainder. The general cumulant
transfer principle is credited to the classical theory.

\section{Analytic derivatives and the tilt parameter}\label{sec:derivatives}

It suffices to analyze one centered, variance-one variable satisfying
\eqref{eq:importedcumulants}. Replacing $S$ by $-S$ preserves the bounds and
changes its third cumulant to $-\kappa_3(S)$. For the next sections write
$K=K_1$ and omit the sign subscripts.

The analytic logarithm has the expansion
\begin{equation}\label{eq:series}
 K(z)=\frac{z^2}{2}+\sum_{k\ge3}c_kz^k,
 \qquad |c_k|\le2(4\delta)^{k-2}.
\end{equation}
The series converges absolutely for $|z|<1/(4\delta)$. Termwise
differentiation and the identity
$\sum_{m\ge0}\binom{m+j}{j}r^m=(1-r)^{-j-1}$ give, for $j\ge3$,
\begin{equation}\label{eq:derivativegeneral}
 |K^{(j)}(z)|
 \le\frac{2j!(4\delta)^{j-2}}{(1-4\delta|z|)^{j+1}},
 \qquad |z|<\frac1{4\delta}.
\end{equation}
In particular,
\begin{equation}\label{eq:derivative34}
 |K'''(z)|\le\frac{48\delta}{(1-4\delta|z|)^4},\qquad
 |K''''(z)|\le\frac{768\delta^2}{(1-4\delta|z|)^5}.
\end{equation}

\begin{lemma}[Real tilt parameters]\label{lem:parameters}
Put $H=(1024\delta)^{-1}$. For real $|h|\le H$,
\begin{equation}\label{eq:realderivatives}
 |K'''(h)|\le50\delta,
 \qquad |K''''(h)|\le800\delta^2,
 \qquad \frac{19}{20}\le K''(h)\le\frac{21}{20}.
\end{equation}
For $0\le x\le(2048\delta)^{-1}$ there is a unique $h\in[0,H]$
with $K'(h)=x$. With $v=K''(h)$ and $a=h\sqrt v$, it satisfies
\begin{align}
 \frac{20}{21}x\le h&\le\frac{20}{19}x,
 &\frac9{10}x\le a&\le\frac{11}{10}x,
 \label{eq:parametercomparisons}\\
 |x-h|&\le25\delta h^2,
 &|a-x|&\le60\delta x^2.
 \label{eq:parameterdifferences}
\end{align}
\end{lemma}
\begin{proof}
For $|h|\le H$, one has $4\delta|h|\le1/256$. The two derivative
bounds in \eqref{eq:realderivatives} follow from
\eqref{eq:derivative34}, since
$48(256/255)^4<50$ and $768(256/255)^5<800$.
Integrating the first and using $K''(0)=1$ gives
\[
 |K''(h)-1|\le50\delta|h|\le\frac{50}{1024}<\frac1{20}.
\]
Thus $K'$ is strictly increasing on $[-H,H]$, with $K'(0)=0$ and
$K'(H)\ge19H/20>(2048\delta)^{-1}$. This proves existence and uniqueness
of $h$. Integrating the variance bounds proves the bounds on $h$ in
\eqref{eq:parametercomparisons}; the bounds on $a$ follow by multiplying by
$\sqrt v\in[\sqrt{19/20},\sqrt{21/20}]$.

Integrating $|K''(u)-1|\le50\delta u$ for $0\le u\le h$ gives
$|x-h|\le25\delta h^2$. Also
\[
 |\sqrt v-1|=\frac{|v-1|}{1+\sqrt v}
 \le\frac{50\delta h}{1+\sqrt{19/20}}<27\delta h.
\]
Hence $|a-x|\le52\delta h^2\le52(20/19)^2\delta x^2<60\delta x^2$.
\end{proof}

\section{A uniform normal approximation under exponential tilting}
\label{sec:tiltedCLT}

For real $|h|\le H$, define the tilted probability measure by
\begin{equation}\label{eq:tilt}
 \frac{d\Pp_h}{d\Pp}=e^{hS-K(h)}.
\end{equation}
Its mean and variance for $S$ are $K'(h)$ and $v=K''(h)$ respectively.
Let
\[
 U_h=\frac{S-K'(h)}{\sqrt v}
\]
be considered under $\Pp_h$.

\begin{lemma}[Tilted Kolmogorov bound]\label{lem:tiltedCLT}
For every real $|h|\le(1024\delta)^{-1}$,
\begin{equation}\label{eq:tiltedBE}
 \dK(\mathcal L_{\Pp_h}(U_h),\Normal)
 \le\left(\frac{52}{\sqrt\pi}
       +\frac{1248}{\pi\sqrt{2\pi}}\right)\delta
 <200\delta.
\end{equation}
\end{lemma}
\begin{proof}
The centered tilted logarithm is
\[
 G_h(z)=K(h+z/\sqrt v)-K(h)-zK'(h)/\sqrt v.
\]
For $T=(52\delta)^{-1}$ and $|s|\le T$, the segment from $h$ to
$h+is/\sqrt v$ lies in $|z|\le(48\delta)^{-1}$, because
\[
 \frac1{1024}+\frac1{52}\sqrt{\frac{20}{19}}<\frac1{48}.
\]
On this segment, \eqref{eq:derivative34} and $v\ge19/20$ give
\[
 \frac{|K'''(z)|}{v^{3/2}}
 \le\frac{48(12/11)^4}{(19/20)^{3/2}}\delta<78\delta.
\]
Taylor's formula with its integral remainder therefore gives
\[
 G_h(is)=-s^2/2+R_h(s),\qquad
 |R_h(s)|\le13\delta|s|^3\le s^2/4.
\]
For the characteristic function $\phi_h$ of $U_h$ under $\Pp_h$,
\begin{equation}\label{eq:tiltedCF}
 |\phi_h(s)-e^{-s^2/2}|
 \le13\delta|s|^3e^{-s^2/4},\qquad |s|\le T.
\end{equation}
The normal-reference form of Esseen's smoothing inequality is
\[
 \dK(\mathcal L_{\Pp_h}(U_h),\Normal)
 \le\frac1\pi\int_{-T}^{T}
          \frac{|\phi_h(s)-e^{-s^2/2}|}{|s|}\,ds
       +\frac{24}{\pi\sqrt{2\pi}\,T};
\]
see \cite[equation~(2.16)]{LPRS}. Since
$\int_\R s^2e^{-s^2/4}\,ds=4\sqrt\pi$, substitution proves
\eqref{eq:tiltedBE}. The displayed constant is $187.817974\ldots<200$.
\end{proof}

\section{Undoing the tilt, including atoms}\label{sec:undo}

We use two elementary facts about the half-line normal Laplace transform
$J$ from \eqref{eq:Jdef}.

\begin{lemma}[Normal integral estimates]\label{lem:J}
For $a\ge0$,
\begin{equation}\label{eq:Jintegral}
 J(a)=\int_0^\infty e^{-ay}\phi(y)\,dy
 \ge\frac{1-e^{-1}}{\sqrt{2\pi}(1+a)}
 >\frac1{4(1+a)}.
\end{equation}
Moreover,
\begin{equation}\label{eq:Jderivative}
 0\le-\frac{d}{da}\log J(a)\le\frac2{1+a}.
\end{equation}
If a real variable $U$ obeys $\dK(U,\Normal)\le\eta$, then for either
$\diamond\in\{\ge,>\}$,
\begin{equation}\label{eq:laplaceerror}
 \left|\E[e^{-aU}\ind_{\{U\diamond0\}}]-J(a)\right|\le2\eta,
 \qquad \Pp(U=0)\le2\eta.
\end{equation}
\end{lemma}
\begin{proof}
Completing the square proves the integral expression for $J$. On $0\le y\le1$,
$e^{-y^2/2}\ge e^{-y}$, so
\[
 J(a)\ge\frac1{\sqrt{2\pi}}\int_0^1e^{-(a+1)y}\,dy
 \ge\frac{1-e^{-1}}{\sqrt{2\pi}(1+a)}.
\]
The numerical inequality $(1-e^{-1})/\sqrt{2\pi}>1/4$ gives
\eqref{eq:Jintegral}.

Let $m(a)$ denote the mean of $y$ under the probability density proportional
to $e^{-ay-y^2/2}$ on $[0,\infty)$. Then
$-(\log J)'(a)=m(a)\ge0$. Differentiation gives
$m'(a)=-\Var_a(y)\le0$, and $m(0)=\sqrt{2/\pi}<1$.
Integration of the derivative of $y e^{-ay-y^2/2}$ gives
$1=a m(a)+\E_a[y^2]$, so $m(a)\le1/a$ for $a>0$.
The two bounds imply $m(a)\le\min(1,1/a)\le2/(1+a)$, proving
\eqref{eq:Jderivative}.

For $a>0$, the identity
\[
 \E[e^{-aU}\ind_{\{U\diamond0\}}]
 =\Pp(U\diamond0)
   -\int_0^\infty a e^{-ay}\Pp(U>y)\,dy
\]
holds for both endpoint conventions. Each of the two terms differs from
its normal counterpart by at most $\eta$. The same assertion at $a=0$
follows directly from the distribution functions and their left limits.
Finally, a jump of a distribution function within Kolmogorov distance
$\eta$ of a continuous function has size at most $2\eta$. This proves
\eqref{eq:laplaceerror} and the atom bound.
\end{proof}

Fix $x$ in \eqref{eq:mainrange} and choose $h$ as in
Lemma~\ref{lem:parameters}. Under $\Pp_h$, $S=x+\sqrt v\,U_h$.
Changing measure gives the exact identity
\begin{equation}\label{eq:undoexact}
 \Pp(S\diamond x)
 =e^{K(h)-hx}\E_h[e^{-aU_h}\ind_{\{U_h\diamond0\}}],
 \qquad a=h\sqrt v.
\end{equation}
Lemmas~\ref{lem:tiltedCLT} and~\ref{lem:J} show that the expectation on
the right differs from $J(a)$ by at most $400\delta$. Dividing by
\eqref{eq:Jintegral} and using $a\le11x/10$ yields
\begin{equation}\label{eq:etabound}
 \Pp(S\diamond x)=e^{-I(x)}J(a)(1+\eta_x),
 \qquad |\eta_x|\le1600\delta(1+a)
 \le1760\delta(1+x)<1800\delta(1+x).
\end{equation}
This proves \eqref{eq:saddlebound}.

For the full correction, use
$\Phibar(x)=e^{-x^2/2}J(x)$ to obtain
\begin{equation}\label{eq:fullidentity}
 \frac{\Pp(S\diamond x)}{\Phibar(x)e^{\Lambda(x)}}
 =\frac{J(a)}{J(x)}(1+\eta_x).
\end{equation}
By \eqref{eq:Jderivative}, \eqref{eq:parametercomparisons}, and
\eqref{eq:parameterdifferences},
\begin{align}
 \left|\log\frac{J(a)}{J(x)}\right|
 &\le\frac{2|a-x|}{1+\min(a,x)}\notag\\
 &\le\frac{180\delta x^2}{1+x}\le180\delta x.
 \label{eq:Jratio}
\end{align}
In the range \eqref{eq:mainrange}, $180\delta x\le180/2048<0.09$.
It follows that $J(a)/J(x)\le1.1$ and
$|J(a)/J(x)-1|\le200\delta x$. Thus \eqref{eq:fullidentity} gives
\[
 \left|\frac{\Pp(S\diamond x)}{\Phibar(x)e^{\Lambda(x)}}-1\right|
 \le200\delta x+1.1\cdot1800\delta(1+x)
 \le2500\delta(1+x),
\]
which is \eqref{eq:fullcorrection}.

\begin{proposition}[Size of a possible endpoint atom]\label{prop:atoms}
In the range \eqref{eq:mainrange},
\begin{equation}\label{eq:atomreference}
 \frac{\Pp(S=x)}{e^{-I(x)}J(a)}\le1800\delta(1+x).
\end{equation}
In the smaller range \eqref{eq:smallrange},
\begin{equation}\label{eq:atomrelative}
 \frac{\Pp(S=x)}{\Pp(S\ge x)}\le3600\delta(1+x).
\end{equation}
The corresponding assertions hold after replacing $S$ by $-S$.
\end{proposition}
\begin{proof}
Changing measure at the point $S=x$ gives
$\Pp(S=x)=e^{-I(x)}\Pp_h(U_h=0)$. The atom estimate in
\eqref{eq:laplaceerror} bounds the last probability by $400\delta$.
The same calculation as \eqref{eq:etabound} proves
\eqref{eq:atomreference}. Under \eqref{eq:smallrange},
$1800\delta(1+x)\le3600/16384<1/2$. Hence
\eqref{eq:saddlebound} bounds $\Pp(S\ge x)$ below by
$\tfrac12e^{-I(x)}J(a)$, yielding \eqref{eq:atomrelative}.
\end{proof}

Thus the treatment of strict and non-strict tails is quantitative. The
argument does not assume a density, a nonlattice distribution, or any fixed
lattice span.

\section{The Legendre correction and the relative-tail windows}
\label{sec:corrections}

The inverse function $h(x)$ is smooth in the range under consideration,
and $h'(x)=1/v$. Differentiating
$I(x)=h(x)x-K(h(x))$ gives $I'(x)=h(x)$. In particular,
\begin{equation}\label{eq:Lambda1}
 \Lambda'(x)=x-h(x),\qquad \Lambda(0)=0.
\end{equation}
Using Lemma~\ref{lem:parameters},
\[
 |x-h(x)|\le25\delta h(x)^2
 \le25(20/19)^2\delta x^2<28\delta x^2.
\]
Integration of \eqref{eq:Lambda1} proves
$|\Lambda(x)|\le(28/3)\delta x^3<10\delta x^3$.

Two further differentiations give
\begin{align}
 \Lambda'''(x)&=\frac{K'''(h)}{v^3},
 \label{eq:Lambda3}\\
 \Lambda''''(x)&=\frac{K''''(h)}{v^4}
                    -\frac{3K'''(h)^2}{v^5}.
 \label{eq:Lambda4}
\end{align}
At zero, $\Lambda(0)=\Lambda'(0)=\Lambda''(0)=0$ and
$\Lambda'''(0)=\kappa_3(S)$. The real derivative and variance bounds yield
\[
 |\Lambda''''(x)|
 \le\left\{800(20/19)^4+7500(20/19)^5\right\}\delta^2
 <11000\delta^2.
\]
Taylor's formula therefore gives
\[
 \left|\Lambda(x)-\frac{\kappa_3(S)x^3}{6}\right|
 \le\frac{11000}{24}\delta^2x^4<500\delta^2x^4.
\]
This proves \eqref{eq:Lambdabound} and \eqref{eq:cubicLambda}. Applying the
whole argument to $\epsilon S$ proves Theorem~\ref{thm:tails} for both
signs, since $\kappa_3(\epsilon S)=\epsilon\kappa_3(S)$.

\begin{proof}[Proof of Corollary~\ref{cor:logbounds}]
Under \eqref{eq:smallrange},
$2500\delta(1+x)\le5000/16384<1/2$. The inequality
$|\log(1+u)|\le2|u|$ for real $|u|\le1/2$ applied to
\eqref{eq:fullcorrection} proves \eqref{eq:logfull}. Combining it with
\eqref{eq:Lambdabound} gives
\[
 \left|\log\frac{p_{\epsilon,\diamond}(x)}{\Phibar(x)}\right|
 \le5000\delta(1+x)+10\delta x^3
 \le12000\delta(1+x^3),
\]
using $1+x\le2(1+x^3)$ for $x\ge0$. Combining
\eqref{eq:logfull} with \eqref{eq:cubicLambda} proves
\eqref{eq:logcubic}.
\end{proof}

The uniform conclusions \eqref{eq:gaussianwindow} and
\eqref{eq:cubicwindow} follow by taking suprema in
\eqref{eq:loggaussian} and \eqref{eq:logcubic}; their upper endpoints
eventually satisfy \eqref{eq:smallrange}. The full-correction conclusion
\eqref{eq:fullwindow} follows directly from \eqref{eq:fullcorrection}.
In terms of $V$ and $W$, the first two standardized windows are
\[
 x=o\bigl((V/W^2)^{1/6}\bigr),\qquad
 x=o\bigl((V/W^2)^{1/4}\bigr),
\]
respectively. They describe relative probabilities, strengthening a
logarithmic moderate-deviation principle.

\section{Sharpness of the Gaussian and cubic-corrected windows}
\label{sec:sharpness}

For each positive integer $m$, take $m$ disjoint singleton batch types,
each sampled with probability $1/m$, and put $t_m=m\log(3/2)$. Each
type then has an independent Poisson arrival count of mean $\log(3/2)$.
Its presence probability is $p=1/3$, so, with $q=1-p=2/3$,
\begin{equation}\label{eq:binomialmodel}
 X_m\ \stackrel{d}{=}\ \Bin(m,p),\qquad
 S_m=\frac{X_m-mp}{\sqrt{mpq}},\qquad
 \delta_m=\frac1{\sqrt{mpq}},\qquad
 \kappa_3(S_m)=(q-p)\delta_m=\frac{\delta_m}{3}.
\end{equation}
These are Poissonized discovery models with $W=B=1$.

For real $x$ with $r=p+pq\delta_m x\in(0,1)$, the exact Legendre transform
of $S_m$ is
\begin{equation}\label{eq:binomialrate}
 I_m(x)=mD(r\Vert p),\qquad
 D(r\Vert p)=r\log\frac rp+(1-r)\log\frac{1-r}{q}.
\end{equation}
To verify this, the cumulant generating function is
\[
 K_m(z)=m\log(q+pe^{\delta_m z})-mp\delta_m z.
\]
The equation $K_m'(h)=x$ makes the tilted success probability equal to $r$,
and substitution in $hx-K_m(h)$ gives \eqref{eq:binomialrate}.

Taylor expansion about the fixed $p\in(0,1)$ gives
\begin{equation}\label{eq:binomialLambda}
 \Lambda_m(x)
 =\frac{(1-2p)\delta_m x^3}{6}
  +\frac{(3pq-1)\delta_m^2x^4}{12}
  +O(\delta_m^3x^5),\qquad \delta_m x\to0.
\end{equation}
For completeness, the derivatives of $D(\,\cdot\,\Vert p)$ at $p$ are
$D''(p)=1/(pq)$, $D'''(p)=(2p-1)/(p^2q^2)$, and
$D''''(p)=2(1-3pq)/(p^3q^3)$; its fifth derivative is bounded in a fixed
neighborhood of $p$. Using $m=(pq\delta_m^2)^{-1}$ in Taylor's formula
proves the stated remainder.

\begin{proposition}[Boundary limits]\label{prop:windows}
Fix $s>0$ and let $\diamond$ denote either $\ge$ or $>$. In the models
\eqref{eq:binomialmodel},
\begin{equation}\label{eq:gaussianboundary}
 \frac{\Pp(S_m\diamond s\delta_m^{-1/3})}
 {\Phibar(s\delta_m^{-1/3})}
 \longrightarrow e^{s^3/18}\ne1.
\end{equation}
Also, with $x_m=s\delta_m^{-1/2}$,
\begin{equation}\label{eq:cubicboundary}
 \frac{\Pp(S_m\diamond x_m)}
 {\Phibar(x_m)\exp\{\kappa_3(S_m)x_m^3/6\}}
 \longrightarrow e^{-s^4/36}\ne1.
\end{equation}
\end{proposition}
\begin{proof}
For either sequence of thresholds, $\delta_m(1+x_m)\to0$, so
Theorem~\ref{thm:tails} gives
\[
 \log\frac{\Pp(S_m\diamond x_m)}{\Phibar(x_m)}
 =\Lambda_m(x_m)+o(1).
\]
In \eqref{eq:binomialLambda}, $p=1/3$ makes the cubic and quartic
coefficients $1/18$ and $-1/36$, respectively. At
$x_m=s\delta_m^{-1/3}$ the cubic term tends to $s^3/18$ and all later
terms vanish. At $x_m=s\delta_m^{-1/2}$, subtracting the cubic term leaves
the quartic limit $-s^4/36$; the remainder is
$O(\delta_m^{1/2})=o(1)$. Exponentiating proves both assertions.
\end{proof}

Thus neither universal relative approximation can in general be extended
to a constant multiple of its stated boundary scale. Individual laws with
cancelling low-order cumulants can have larger windows. The example is
lattice-valued, but no threshold rounding is required: the theorem covers
all real $x$ and both endpoint conventions, and the controlled endpoint
atom is negligible at these scales.

\section{Uniform fixed-sample mean and variance corrections}
\label{sec:uniform-correction}

The fixed-law correction identities in \cite[Section~6]{LongFoundation}
admit absolute error bounds uniform in the sampling law. This remains true
when the law and the weights vary with the sample size. The result identifies the appropriate variance after
removing the Poisson sample-size fluctuation; it does not by itself provide
a fixed-sample Berry--Esseen theorem.

For ordinary unweighted singleton occupancy, uniform $O(1)$ mean and
residual-variance corrections are already given by
Hwang and Janson \cite[Theorem~2.3]{HJ}. Their Theorem~2.1 and
Corollary~2.5 also supply a uniform local approximation and a changing-law
CLT in that singleton model. The result below treats overlapping batches
and nonnegative weights with explicit batch-scale constants. It does not
extend their local limit theorem to the batch model.

Let $A_1,A_2,\ldots$ be independent, identically distributed random finite
subsets of a countable set $\mathcal S$. Let $N$ be an independent rate-one
Poisson process. Assign deterministic weights $h_i\ge0$ and assume
\[
 T(A_1):=\sum_{i\in A_1}h_i\le W\quad\hbox{almost surely},
 \qquad 0<W<\infty.
\]
No uniform bound on batch cardinalities is required. Define
\begin{align*}
 D_{h,n}&=\sum_{i\in\bigcup_{k=1}^n A_k}h_i,
 &X_h(t)&=D_{h,N(t)},\\
 b_{h,n}&=\E D_{h,n},
 &m_h(t)&=\E X_h(t),\\
 s_{h,n}^2&=\Var(D_{h,n}),
 &V_h(t)&=\Var(X_h(t)).
\end{align*}
Set
\[
 \lambda_i=\Pp(i\in A_1),\qquad
 \lambda_{ij}=\Pp(i,j\in A_1),\qquad \lambda_{ii}=\lambda_i,
\]
and
\[
 f_h(t)=\sum_i h_i\lambda_i e^{-t\lambda_i},\qquad
 \mu_1=\E T(A_1),\qquad \mu_2=\E T(A_1)^2.
\]
Species of weight zero or zero arrival rate can be omitted. In particular,
$f_h$ is finite and $m_h'(t)=f_h(t)$, by domination by
$\sum_i h_i\lambda_i=\mu_1\le W$.
All sums over $(i,j)$ below are over ordered pairs, including the diagonal.

\begin{theorem}[Uniform correction with the actual residual variance]
\label{thm:uniform-correction}
For every integer $n\ge1$, put
\begin{equation}\label{eq:tau-definition}
 \tau_{h,n}^2=V_h(n)-n f_h(n)^2.
\end{equation}
Then
\begin{equation}\label{eq:uniform-mean-correction}
 0\le b_{h,n}-m_h(n)\le2\mu_1\le2W
\end{equation}
and
\begin{equation}\label{eq:uniform-variance-correction}
 \left|s_{h,n}^2-\tau_{h,n}^2\right|
 \le2\mu_2+(2+2e^{-2})\mu_1^2
 \le(4+2e^{-2})W^2.
\end{equation}
More precisely,
\begin{equation}\label{eq:uniform-asymmetric-correction}
 -2(\mu_2+\mu_1^2)-2e^{-2}\mu_1^2
 \le s_{h,n}^2-\tau_{h,n}^2
 \le2(\mu_2+\mu_1^2).
\end{equation}
The quantity $\tau_{h,n}^2$ is nonnegative. It has the interpretations
\begin{equation}\label{eq:tau-projection}
 \tau_{h,n}^2
 =\Var\bigl(X_h(n)-f_h(n)N(n)\bigr)
 =\min_{a\in\mathbb R}\Var\bigl(X_h(n)-aN(n)\bigr)
\end{equation}
and
\begin{align}
 \tau_{h,n}^2
 &=n\Var\!\left(\sum_{i\in A_1}h_i e^{-n\lambda_i}\right)\notag\\
 &\quad+\sum_{i,j}h_i h_j e^{-n(\lambda_i+\lambda_j)}
       \bigl(e^{n\lambda_{ij}}-1-n\lambda_{ij}\bigr).
 \label{eq:tau-positive}
\end{align}
Consequently, along any sequence of models with
$s_{h,n}^2/W^2\to\infty$ or $\tau_{h,n}^2/W^2\to\infty$,
\[
 s_{h,n}^2\sim\tau_{h,n}^2.
\]
The same statements hold in the unweighted model with $h_i=1$ and $W=B$.
\end{theorem}

\begin{proof}
First suppose only finitely many species are retained. The countable
passage is given below. For $n\ge1$, write
\[
 \delta_n(x)=e^{-nx}-(1-x)^n,\qquad 0\le x\le1.
\]
We will use
\begin{equation}\label{eq:uniform-delta-lemma}
 \delta_n(x)\ge0,\qquad
 |\delta_n'(x)|\le2,\qquad
 \delta_n(x)\le2x.
\end{equation}
Nonnegativity follows from $1-x\le e^{-x}$. For $n=1$ the derivative bound
is immediate. If $n\ge2$ and $0\le x\le1/2$, inserting
$e^{-(n-1)x}$ between $(1-x)^{n-1}$ and $e^{-nx}$ gives
\begin{align*}
 |\delta_n'(x)|
 &\le n(n-1)x^2e^{-(n-1)x}+nxe^{-(n-1)x}\\
 &\le2(4e^{-2}+e^{-1})<2.
\end{align*}
Here we used $0\le-\log(1-x)-x\le x^2$ and the bounds
$ye^{-y}\le e^{-1}$ and $y^2e^{-y}\le4e^{-2}$ for $y\ge0$.
For $1/2<x\le1$,
\[
 |\delta_n'(x)|
 \le n2^{-(n-1)}+ne^{-n/2}\le1+2/e<2.
\]
The derivatives at endpoints have the usual one-sided meanings. Since
$\delta_n(0)=0$, the final inequality in
\eqref{eq:uniform-delta-lemma} follows by integration.

The exact mean formulas now give
\[
 b_{h,n}-m_h(n)=\sum_i h_i\delta_n(\lambda_i),
\]
which proves \eqref{eq:uniform-mean-correction}.

For an ordered pair $(i,j)$, abbreviate
\[
 a=\lambda_i,\quad b=\lambda_j,\quad c=\lambda_{ij},
 \qquad u=a+b-c,\quad w=a+b-ab.
\]
Both $u$ and $w$ belong to $[0,1]$ and are at least $\max(a,b)$.
These assertions include $i=j$ and all endpoint values of the rates.
The exact fixed-sample covariance of the two discovery indicators is
\[
 (1-u)^n-(1-a)^n(1-b)^n=(1-u)^n-(1-w)^n.
\]
The corresponding Poissonized covariance at time $n$ is
$e^{-nu}-e^{-n(a+b)}$. Thus
\begin{equation}\label{eq:uniform-covsplit}
 s_{h,n}^2=V_h(n)+E_{h,n}-R_{h,n},
\end{equation}
where
\begin{align*}
 E_{h,n}
 &=\sum_{i,j}h_i h_j\{\delta_n(w)-\delta_n(u)\},\\
 R_{h,n}
 &=\sum_{i,j}h_i h_j e^{-n(a+b)}(e^{nab}-1).
\end{align*}
The Lipschitz bound gives
\[
 |E_{h,n}|
 \le2\sum_{i,j}h_i h_j(c+ab)
 =2(\mu_2+\mu_1^2).
\]
The linear term in $R_{h,n}$ is precisely $nf_h(n)^2$. Using
$e^z-1-z\le z^2e^z/2$ for $z\ge0$, together with
$w=a+b-ab\ge(a+b)/2$, yields
\begin{align*}
 0\le R_{h,n}-nf_h(n)^2
 &\le\frac{n^2}{2}\sum_{i,j}h_i h_j a^2b^2e^{-nw}\\
 &\le\frac12
    \left(\sum_i h_i\lambda_i(n\lambda_i)
                              e^{-n\lambda_i/2}\right)^2\\
 &\le2e^{-2}\mu_1^2.
\end{align*}
The last step uses $ye^{-y/2}\le2/e$ for $y\ge0$.
Combining this with \eqref{eq:uniform-covsplit} proves
\eqref{eq:uniform-asymmetric-correction} and hence
\eqref{eq:uniform-variance-correction}.

For the projection interpretation, the elementary Poisson identity
$\E[N(n)g(N(n))]=n\E[g(N(n)+1)]$, applied after averaging the batches,
gives
\begin{align*}
 \operatorname{Cov}(X_h(n),N(n))
 &=n\E[D_{h,N(n)+1}-D_{h,N(n)}]\\
 &=n\sum_i h_i\lambda_i e^{-n\lambda_i}=nf_h(n).
\end{align*}
Since $\Var(N(n))=n$, completion of the square proves
\eqref{eq:tau-projection}. Alternatively, separate the linear term in
\[
 V_h(n)=\sum_{i,j}h_i h_j e^{-n(\lambda_i+\lambda_j)}
                                  (e^{n\lambda_{ij}}-1).
\]
That linear term is
$n\E[(\sum_{i\in A_1}h_i e^{-n\lambda_i})^2]$.
Subtracting $nf_h(n)^2$ gives \eqref{eq:tau-positive}, whose summands are
nonnegative.

For countably many species, take an increasing finite exhaustion $J_k$ and
intersect every batch with $J_k$. Empty intersections cause no contribution;
identical intersections may be coalesced if a finite type space is desired.
The same total-weight bound holds. The truncated fixed-sample counts
increase to $D_{h,n}$ and are bounded by $nW$, while their Poissonized
counterparts increase to $X_h(n)$ and are bounded by $WN(n)$.
These bounds imply $L^2$ convergence, and hence convergence of all means
and variances in the theorem. The sums defining $f_h(n)$, $\mu_1$, and
$\mu_2$ converge by monotone convergence; the corresponding weighted
batch sums are bounded by $W$. Thus the finite inequalities and the
projection identity pass to the limit. The nonnegative series in
\eqref{eq:tau-positive} also pass by monotone convergence.

For completeness, the signed discrepancy sums are absolutely convergent:
\[
 \sum_{i,j}h_i h_j|\delta_n(w)-\delta_n(u)|
 \le2(\mu_2+\mu_1^2)<\infty,
\]
and $R_{h,n}\le n\mu_1^2$, because
$e^{-nw}-e^{-n(a+b)}\le nab$. The fixed-sample covariance series is
absolutely bounded by $n(\mu_2+\mu_1^2)$ using the Lipschitz constant $n$
of $x\mapsto(1-x)^n$. Thus the displayed sum manipulations are also
justified directly. The final asymptotic assertion follows by dividing
\eqref{eq:uniform-variance-correction} by the variance scale that diverges
relative to $W^2$.
\end{proof}

\begin{remark}[A moment-only version of the correction]
\label{rem:uniform-moment-correction}
The inequalities retaining $\mu_1$ and $\mu_2$, and identities
\eqref{eq:tau-projection}--\eqref{eq:tau-positive}, require only
$\E T(A_1)^2<\infty$, without an almost-sure upper bound $W$.
Indeed, use domination by $\sum_{k=1}^nT(A_k)$ for fixed sample size and
by $\sum_{k=1}^{N(n)}T(A_k)$ after Poissonization. Their second moments are
finite, and the same incidence sums and discrepancy estimates apply.
This observation concerns mean and variance corrections only; it does not
assert a normal-approximation theorem for unbounded batches or weights.
\end{remark}

\begin{remark}[Why the Poissonized variance cannot simply be substituted]
\label{rem:uniform-variance-cancellation}
Take $B=1$ and a uniform singleton law on
$M_n=\lceil n^{3/2}\rceil$ species. With $a_n=n/M_n$,
\[
 V(n)=M_n e^{-a_n}(1-e^{-a_n})\sim n,
 \qquad f(n)=e^{-a_n},
\]
whereas
\[
 \tau_n^2
 =M_n e^{-2a_n}(e^{a_n}-1-a_n)
 \sim\frac{n^2}{2M_n}\sim\frac{\sqrt n}{2}.
\]
The uniform correction therefore proves
$s_n^2\sim\sqrt n/2$, so both variances diverge but
$V(n)/s_n^2\sim2\sqrt n\to\infty$. Thus the fixed-law variance
equivalence from qualitative de-Poissonization is not uniform over changing
laws. This example does not disprove a fixed-sample bound $CB/s_n$.
The corrected variance estimate supplies the right scale, but a uniform
normal approximation at that scale, or a suitable quantitative conditional
limit theorem, remains an additional proof obligation.
\end{remark}

\begin{remark}[The uniform remainder orders cannot be improved]
\label{rem:uniform-correction-order}
For the uniform singleton law on $n$ species with $n$ draws, write
$u_n=(1-1/n)^n$ and $v_n=(1-2/n)^n$ (for $n\ge2$). Then
\[
 b_n-m(n)=n(e^{-1}-u_n),
 \qquad
 s_n^2=nu_n(1-u_n)+n(n-1)(v_n-u_n^2),
\]
and $\tau_n^2=n(e^{-1}-2e^{-2})$. The elementary expansions
\begin{align*}
 u_n&=e^{-1}\left(1-\frac1{2n}-\frac5{24n^2}+O(n^{-3})\right),\\
 u_n^2&=e^{-2}\left(1-\frac1n-\frac1{6n^2}+O(n^{-3})\right),\\
 v_n&=e^{-2}\left(1-\frac2n-\frac2{3n^2}+O(n^{-3})\right)
\end{align*}
give
\[
 b_n-m(n)\longrightarrow\frac1{2e},\qquad
 s_n^2-\tau_n^2\longrightarrow
       \frac{3e^{-2}-e^{-1}}2>0.
\]
Replacing every singleton by a disjoint block of size $B$ multiplies the
mean discrepancy by $B$ and the variance discrepancy by $B^2$. Equivalently,
assigning weight $W$ to every singleton multiplies them by $W$ and $W^2$.
Thus the uniform remainder scales $O(W)$ and $O(W^2)$ cannot be replaced
by $o(W)$ and $o(W^2)$ over changing laws, even when the corrected and
fixed-sample variances diverge. This concerns the orders of the correction
errors, not the optimal constants in
\eqref{eq:uniform-mean-correction}--\eqref{eq:uniform-variance-correction}.
\end{remark}

\section{A sharper envelope for multiple coverage}
\label{sec:sharper-threshold}

Let the batch law be $\mu$, let $P_C$ denote the independent Poisson clock
of batch type $C$, with rate $\mu(C)$, and put
\[
 \lambda_i=\sum_{C\ni i}\mu(C),\qquad
 \lambda_{ij}=\sum_{C\ni i,j}\mu(C),\qquad
 N_i(t)=\sum_{C\ni i}P_C(t).
\]
We allow nonnegative species weights $w_i$ and assume that every batch of
positive probability satisfies
\begin{equation}\label{eq:threshold-weight-budget}
 \sum_{i\in C}w_i\le W,\qquad 0<W<\infty.
\end{equation}
Species of zero rate or zero weight may be discarded. Let $r_i\in\{1,\ldots,r\}$,
where $r$ is a positive integer, and consider
\[
 Y_r(t)=\sum_i w_i\boldsymbol 1_{\{N_i(t)\ge r_i\}},\qquad
 m_r(t)=\mathbb EY_r(t),\qquad V_r(t)=\operatorname{Var}Y_r(t).
\]
The unweighted model has $w_i=1$ and $W=B$.
The bound $Y_r(t)\le WN(t)$ justifies finite-species approximation and all
moment limits used below.

Define
\begin{equation}\label{eq:mode-envelope-constant}
 \ell_r=1+\sum_{k=1}^{r-1}\frac{e^{-k}k^k}{k!},\qquad
 \widetilde W_r=W\sqrt{\ell_r}.
\end{equation}
In particular, $\ell_1=1$ and $\ell_2=1+e^{-1}$.

\begin{lemma}[Integrated Poisson mode envelope]
\label{lem:poisson-mode-envelope}
For $p_k(x)=e^{-x}x^k/k!$ and
\[
 g_r(x)=\max_{0\le k\le r-1}p_k(x),\qquad x\ge0,
\]
one has
\begin{equation}\label{eq:integrated-mode-envelope}
 \int_0^\infty g_r(x)\,dx=\ell_r,
 \qquad
 \ell_r\le1+\sqrt{\frac{2(r-1)}\pi},
 \qquad
 \ell_r\sim\sqrt{\frac{2r}\pi}\quad(r\to\infty).
\end{equation}
Moreover, for all $a,b>0$,
\begin{equation}\label{eq:two-rate-mode-envelope}
 \int_0^\infty g_r(au)g_r(bu)\,du
 \le\frac{\ell_r}{\max\{a,b\}},
\end{equation}
and $\ell_r$ is the best possible constant in this last inequality when
all positive $a,b$ are allowed.
\end{lemma}

\begin{proof}
The ratio $p_k(x)/p_{k-1}(x)=x/k$ identifies the mode. The maximizing
index is $\lfloor x\rfloor$ until it reaches $r-1$, and is $r-1$
thereafter; at integral transition points there can be two equal modes.
For $r\ge2$, the functions $g_r$ and $g_{r-1}$ agree on $[0,r-1]$, while
on $[r-1,\infty)$ their difference is $p_{r-1}-p_{r-2}$.
Since $p'_{r-1}=p_{r-2}-p_{r-1}$,
\[
 \int_0^\infty(g_r-g_{r-1})(x)\,dx
 =p_{r-1}(r-1).
\]
Together with $\int_0^\infty g_1(x)\,dx=1$, this proves the exact formula.
Stirling's lower bound $k!\ge\sqrt{2\pi k}(k/e)^k$, for $k\ge1$, gives
\[
 \ell_r\le1+\frac1{\sqrt{2\pi}}\sum_{k=1}^{r-1}k^{-1/2}
 \le1+\sqrt{\frac{2(r-1)}\pi}.
\]
Stirling's asymptotic formula and summation of $k^{-1/2}$ give the stated
asymptotic.

For \eqref{eq:two-rate-mode-envelope}, keep the factor with the larger
rate and bound the other by one. To see sharpness for this particular
envelope inequality, take $a=1$ and let $b\downarrow0$.
Then $g_r(bu)\to1$ for every $u$, and domination by the integrable
function $g_r(u)$ gives convergence of the integral to $\ell_r$.
This is a sharpness statement about the deterministic envelope, not about
normal approximation for the stochastic model.
\end{proof}

\begin{theorem}[Weighted threshold bracket bounds]
\label{thm:sharper-threshold-bracket}
For the terminal Doob martingale of $Y_r(t)$, all jumps have absolute
value at most $W$. Its predictable quadratic variation $Q_{r,t}$ obeys
\begin{align}
 \mathbb EQ_{r,t}&=V_r(t),\\
 \operatorname{Var}Q_{r,t}&\le2\ell_rW^2V_r(t),
       \label{eq:sharper-threshold-bracket-variance}\\
 \log\mathbb E e^{\theta(Q_{r,t}-V_r(t))}
 &\le\frac{V_r(t)}{2\ell_rW^2}
       \bigl(e^{2\ell_rW^2\theta}-1-2\ell_rW^2\theta\bigr),
       \qquad \theta\in\mathbb R.
       \label{eq:sharper-threshold-bracket-mgf}
\end{align}
Here and below a zero variance is treated as the constant-variable case.
\end{theorem}

\begin{proof}
Fix the terminal time $t$. The jump in the conditional expectation of
the unweighted contribution of species $i$, when one arrival containing
$i$ occurs at time $s$, is
\[
 h_i(s)=\boldsymbol 1_{\{N_i(s-)<r_i\}}
 p_{r_i-1-N_i(s-)}\bigl(\lambda_i(t-s)\bigr).
\]
Consequently a batch $C$ produces jump
$H_C(s)=\sum_{i\in C}w_i h_i(s)\in[0,W]$, and
\begin{equation}\label{eq:sharper-threshold-bracket-expansion}
 Q_{r,t}=\sum_{i,j}w_iw_j\lambda_{ij}
       \int_0^t h_i(s)h_j(s)\,ds.
\end{equation}
The identity $\mathbb EQ_{r,t}=V_r(t)$ is the martingale isometry.

Pointwise, $h_i(s)\le g_r(\lambda_i(t-s))$.
Lemma~\ref{lem:poisson-mode-envelope} therefore gives
\[
 0\le\int_0^t h_i(s)h_j(s)\,ds
 \le\frac{\ell_r}{\max\{\lambda_i,\lambda_j\}}.
\]
Write the bracket as a positive sum $Q_{r,t}=\sum_{i,j}a_{ij}F_{ij}$,
where
\[
 a_{ij}=\ell_r w_iw_j\frac{\lambda_{ij}}
                            {\max\{\lambda_i,\lambda_j\}},\qquad
 F_{ij}=\frac{\max\{\lambda_i,\lambda_j\}}{\ell_r}
              \int_0^t h_i(s)h_j(s)\,ds\in[0,1].
\]
Terms with $a_{ij}=0$ are omitted. Each $F_{ij}$ is a function of the
independent batch paths meeting $\{i,j\}$. For a fixed batch $C$, its
coefficient load is at most
\begin{align*}
 \sum_{i\in C\ \mathrm{or}\ j\in C}a_{ij}
 &\le2\ell_r\sum_{i\in C}\frac{w_i}{\lambda_i}
                         \sum_jw_j\lambda_{ij}\\
 &\le2\ell_r W\sum_{i\in C}w_i
 \le2\ell_rW^2.
\end{align*}
The middle inequality follows from
\[
 \sum_jw_j\lambda_{ij}
 =\sum_{D\ni i}\mu(D)\sum_{j\in D}w_j\le W\lambda_i.
\]
The weighted variance-drop and exponential weighted-overlap inequalities
from \cite{LongFoundation}, applied to this positive decomposition with
load bound $2\ell_rW^2$, prove
\eqref{eq:sharper-threshold-bracket-variance}--\eqref{eq:sharper-threshold-bracket-mgf}.
The countable case follows by nonnegative partial sums. In particular,
$Q_{r,t}\le W^2t$ supplies direct domination for every required bracket
moment and exponential moment.
\end{proof}

\begin{corollary}[Normal approximation with growing thresholds]
\label{cor:sharper-threshold-normal}
If $V_r(t)>0$, put
\[
 \eta_r=\frac{\ell_rW^2}{V_r(t)}.
\]
There is an absolute constant $C$ such that
\begin{align}
 d_K\!\left(\frac{Y_r(t)-m_r(t)}{\sqrt{V_r(t)}},N(0,1)\right)
 &\le\min\{1,C\eta_r^{1/5}\},
       \label{eq:sharper-threshold-fifth}\\
 d_K\!\left(\frac{Y_r(t)-m_r(t)}{\sqrt{V_r(t)}},N(0,1)\right)
 &\le\min\!\left\{1,
 C\bigl[\eta_r\log(e+\eta_r^{-1})\bigr]^{1/4}\right\}.
       \label{eq:sharper-threshold-log}
\end{align}
Thus, for a sequence of models, weights, thresholds, and times, the
standardized counts converge to $N(0,1)$ whenever
\begin{equation}\label{eq:sharper-threshold-clt-condition}
 \frac{\ell_rW^2}{V_r(t)}\longrightarrow0.
\end{equation}
For growing $r$, this condition is equivalent to
$\sqrt r\,W^2/V_r(t)\to0$.
\end{corollary}

\begin{proof}
Apply the generic martingale normal-approximation arguments of
\cite{LongFoundation} with $\widetilde W_r=W\sqrt{\ell_r}$ in place of
the jump-scale parameter. Indeed, $W\le\widetilde W_r$, while the bracket
has expectation $V_r(t)$ and coefficient load at most
$2\widetilde W_r^2$. These are exactly the bounded-jump, variance, and
exponential-bracket inputs used in the fifth-root and logarithmic
fourth-root estimates. Their constants are universal and do not depend
on $r$. The deterministic-grid justification is unchanged: between
arrivals $|h_i'(s)|\le\lambda_i\le1$, so $|H_C'(s)|\le W$ and the
bracket density has derivative bounded by $2W^2$.
The CLT follows because $\eta\log(e+\eta^{-1})\to0$ as $\eta\downarrow0$.
\end{proof}

\begin{corollary}[Actual-variance threshold concentration]
\label{cor:sharper-threshold-concentration}
For real $\theta$ with $|\theta|<1/(2\widetilde W_r)$,
\begin{equation}\label{eq:sharper-threshold-actual-mgf}
 \log\mathbb E e^{\theta(Y_r(t)-m_r(t))}
 \le\frac{V_r(t)\theta^2}{2(1-2\widetilde W_r|\theta|)}.
\end{equation}
Consequently, for $x>0$,
\begin{equation}\label{eq:sharper-threshold-actual-tail}
 \mathbb P\{|Y_r(t)-m_r(t)|\ge x\}
 \le2\exp\!\left\{-\frac{x^2}{2(V_r(t)+2\widetilde W_rx)}\right\}.
\end{equation}
The linear scale is $2W\sqrt{\ell_r}=O(Wr^{1/4})$, while the
quadratic coefficient remains the actual variance.
\end{corollary}

\begin{proof}
Write $M=Y_r(t)-m_r(t)$, $V=V_r(t)$, and $A=\widetilde W_r$.
Let $\psi_A(u)=(e^{Au}-1-Au)/A^2$ for $u\ge0$.
The compensated-Poisson exponential martingale and $0\le H_C\le A$
give, for $q>1$,
\[
 \mathbb E e^{q\theta M-\psi_A(q|\theta|)Q_{r,t}}\le1.
\]
H\"older's inequality and the uncentered form of
\eqref{eq:sharper-threshold-bracket-mgf} imply
\begin{equation}\label{eq:sharper-threshold-holder}
 \log\mathbb E e^{\theta M}
 \le\frac{q-1}{q}\frac{V}{2A^2}
 \left[\exp\!\left\{\frac{2A^2\psi_A(q|\theta|)}{q-1}\right\}-1\right].
\end{equation}
For $a=A|\theta|\in(0,1/2)$ choose $q=1+a$.
The elementary bounds
\[
 e^u-1-u\le\frac{u^2}{2(1-u/3)}\quad(0\le u<3),\qquad
 e^z-1\le\frac{z}{1-z/2}\quad(0\le z<2)
\]
then bound \eqref{eq:sharper-threshold-holder} by
\[
 \frac{V\theta^2}{2}\,
 \frac{q}{1-qa/3-q^2a/2}
 \le\frac{V\theta^2}{2(1-2a)}.
\]
For completeness, all denominators are positive because
\[
 1-\frac{qa}{3}-\frac{q^2a}{2}-q(1-2a)
 =\frac a6+\frac{2a^2}{3}-\frac{a^3}{2}\ge0.
\]
This also ensures that the second exponential-series bound is applied
with argument less than two. The case $\theta=0$ is immediate.
For $V>0$, Chernoff's inequality at
$|\theta|=x/(V+2Ax)$ gives each one-sided tail bound; summing proves
\eqref{eq:sharper-threshold-actual-tail}.
\end{proof}

\begin{remark}[Scope of the improvement]
The new parameter $\ell_r=\Theta(\sqrt r)$ replaces a parameter of order
$r$ in the earlier bracket estimate. It improves the permitted growth
of the thresholds, the generic normal-approximation bounds, and the
actual-variance Bernstein scale. It does not prove an optimal
$O(W/\sqrt{V_r})$ Kolmogorov bound for thresholds greater than one,
or show that any threshold dependence in such a bound is necessary.
The sharpness assertion in Lemma~\ref{lem:poisson-mode-envelope} concerns
only the displayed deterministic point-mass envelope.
\end{remark}

\section{Scope and remaining questions}\label{sec:scope}

The relative-tail estimates in Sections~\ref{sec:main}--\ref{sec:sharpness}
are uniform across Poissonized first-discovery models subject to
\eqref{eq:batchweight} and $V>0$. This includes varying laws, times, weights,
and countable species spaces, as well as unbounded batch cardinalities
with bounded total weight. The countable passage is already part of the
imported analytic and cumulant theorem; no finite-species approximation
error is added to those estimates. The Gaussian and cubic-corrected
relative-tail windows are sharp for this class.

Section~\ref{sec:uniform-correction} identifies the residual variance after
removing the Poisson sample-size fluctuation and proves a uniform
$O(W^2)$ correction. It also shows why the Poissonized variance alone is
not an adequate substitute when the law varies. A uniform fixed-sample
Kolmogorov bound of order $W/s_{h,n}$ still requires a distributional
argument at the corrected variance scale. Relative-tail transfer requires
additional control of errors at the tail scale, and is not a consequence
of the moment corrections alone.

Section~\ref{sec:sharper-threshold} improves the threshold dependence of
bracket, generic normal-approximation, and concentration estimates. The
factor $\ell_r$ is optimal for the particular deterministic two-rate
envelope, but no necessity of threshold dependence in an optimal
probabilistic bound is asserted. A threshold-uniform Kolmogorov bound of
order $W/\sqrt{V_r}$ remains open here. The first-discovery cumulant theorem
does not directly apply to higher coverage or to signed species weights.

No joint occupancy-spectrum CLT, functional limit theorem, or uniform
fixed-sample moderate-deviation theorem is proved in this note. Local
limits require appropriate arithmetic or smoothing information. These
questions and the unsuccessful proof routes are recorded separately in
the repository research-status note. All manuscripts remain research
drafts; the internal proof audits and deterministic finite calculations
do not constitute independent peer review or formal verification.

\begin{thebibliography}{9}

\bibitem{DJS}
H. D\"oring, S. Jansen, and K. Schubert.
The method of cumulants for the normal approximation.
\emph{Probability Surveys} \textbf{19} (2022), 185--270.
\href{https://doi.org/10.1214/22-PS7}{doi:10.1214/22-PS7};
\href{https://arxiv.org/abs/2102.01459}{arXiv:2102.01459}.
Theorem~2.1 and Section~6.

\bibitem{LPRS}
J. L. Lebowitz, B. Pittel, D. Ruelle, and E. R. Speer.
Central limit theorems, Lee--Yang zeros, and graph-counting polynomials.
\emph{Journal of Combinatorial Theory, Series A} \textbf{141} (2016),
147--183.
\href{https://doi.org/10.1016/j.jcta.2016.02.009}{doi:10.1016/j.jcta.2016.02.009};
\href{https://arxiv.org/abs/1408.4153}{arXiv:1408.4153}.
Equation~(2.16).

\bibitem{HJ}
H.-K. Hwang and S. Janson.
Local limit theorems for finite and infinite urn models.
\emph{The Annals of Probability} \textbf{36} (2008), 992--1022.
\href{https://doi.org/10.1214/07-AOP350}{doi:10.1214/07-AOP350};
\href{https://arxiv.org/abs/math/0604397}{arXiv:math/0604397}.
Theorems~2.1 and 2.3, Corollary~2.5, and Remark~5.1.

\bibitem{LongFoundation}
C. D. Long.
\emph{A fluctuation theory for bounded-batch species discovery}.
Research draft, November 12, 2026; repository revision November 26, 2026.
\href{https://github.com/long-mathematics/bounded-batch-discovery}{Repository},
\path{papers/bounded_batch_discovery.tex}.

\bibitem{LongOptimal}
C. D. Long.
\emph{Optimal Berry--Esseen bounds for bounded-batch discovery}.
Research draft, November 26, 2026.
\href{https://github.com/long-mathematics/bounded-batch-discovery}{Repository},
\path{papers/optimal_berry_esseen.tex}.

\end{thebibliography}
\end{document}
